\documentclass[10pt,a4paper]{amsart}
\usepackage[margin=19mm]{geometry}
\usepackage{amsmath,amssymb,mathtools}
\usepackage[scr=rsfs,cal=euler]{mathalfa}
\usepackage{xfrac}
\usepackage[unicode,hidelinks,bookmarks=false]{hyperref}
\newcommand{\ie}{i.e.\ }
\newcommand{\cf}{cf.\ }
\newcommand{\resp}{respectively\ }
\newcommand{\Subsection}{\subsection}
\providecommand{\nobreakdash}{\nobreak\hbox{-}}
\setlength{\emergencystretch}{2em}
\multlinegap0pt
\usepackage{bm}
\usepackage{bbm}
\usepackage{dsfont}
\usepackage{xspace}
\usepackage{cancel}
\usepackage{mathtools}
\usepackage[shortlabels]{enumitem}
\usepackage{amsthm}
\makeatletter
\@ifundefined{theorem}{\newtheorem{theorem}{Theorem}}{}
\@ifundefined{lemma}{\newtheorem{lemma}[theorem]{Lemma}}{}
\@ifundefined{observation}{\newtheorem{observation}[theorem]{Observation}}{}
\makeatother
\def\H{\mathcal{H}}
\def\nat{\mathbb{N}}
\newcommand{\tLip}{\textnormal{Lip}}
\title[Perturbations of measures and sets having curves in d direct]{Perturbations of measures and sets having curves in d directions}
\author{Jakub Tak\'a\v{c}}
\date{Source: arXiv:2604.15887v1 (CC BY 4.0)}
\begin{document}
\maketitle

\noindent\emph{Source and licence.} Perturbations of measures and sets having curves in d directions. Version arXiv:2604.15887v1 (CC BY 4.0), distributed under the Creative Commons Attribution 4.0 International licence, as recorded in the arXiv licence metadata for this exact version and shown on the abstract page https://arxiv.org/abs/2604.15887v1. Full rights notice: licenses/arxiv-2604.15887-CC-BY-4.0.txt.\par\smallskip

\noindent\textit{Editorial scope.} Counted unit: the Theorem whose source label is \texttt{None} in the article (source0.tex lines 1721--1730, source label \texttt{None}), delivered together with the article's own complete original proof of it (source0.tex lines 1731--1812, one \texttt{proof} environment of 82 lines). No mathematics is written, repaired or re-derived: statement and proof are the article's own text line for line. Required context retained uncounted: O:finiteness-subspaces (lines 1634--1636); T:StV-main (lines 1643--1653); L:semicont (lines 1710--1716). Every definition, display and label of those lines is retained unchanged. Omitted from this package and not redistributed: 1--1633, 1637--1642, 1654--1709, 1717--1720, 1813--2716. Nothing inside a retained excerpt is omitted or altered: every retained body is delivered exactly as the article's lines are written, so no omission receipt of any kind accompanies this package. Citations: the retained excerpts carry no citation command at all, so no bibliography entry of the article is transcribed and no reference is redistributed. Reference adaptation: none was needed and none was made; every reference inside the delivered unit resolves inside it, so no unresolved reference or citation is produced. Printed slips kept literal: no printed word, symbol, hypothesis or number of the counted statement or of its proof was altered. Layout and class: the article's own class is \texttt{article} with packages including inputenc, lmodern, a4wide, amsmath, amssymb, mathrsfs, amsthm, natbib, enumitem, geometry, stmaryrd, authblk, url, hyperref; the wrapper is the lane's pinned \texttt{amsart} editorial wrapper at 10pt on A4, which restates the theorem-like environments the retained text needs and adds the article's own macro definitions that the retained text needs (\texttt{\textbackslash H}, \texttt{\textbackslash nat}, \texttt{\textbackslash tLip}), with extra wrapper packages (bm, bbm, dsfont, xspace, cancel, mathtools, enumitem). The delivered numbering is the unit's own. Assets: no retained span contains a figure environment, a graphics-inclusion command, a PDF-page inclusion, a table or a TikZ picture; no image file is copied into this package, the package declares no asset dependency and no image file is redistributed with it. Licence: version arXiv:2604.15887v1 of this article is distributed under the Creative Commons Attribution 4.0 International licence, as recorded in the arXiv licence metadata for this exact version and shown on the abstract page https://arxiv.org/abs/2604.15887v1. A full rights notice is delivered at licenses/arxiv-2604.15887-CC-BY-4.0.txt. Characters: every retained excerpt is pure ASCII, so the delivered engine, XeLaTeX with its default Latin Modern text font, reports no missing character.
\begin{observation}\label{O:finiteness-subspaces}
    Suppose $V$ is a finite-dimensional Banach space and $\mathcal{F}\subset V$ is a finite union of $d$-dimensional affine subspaces. If $F\subset \mathcal{F}$ is compact, then $\H^d(F)<\infty$.
\end{observation}

\begin{theorem}\label{T:StV-main}
    Let $X$ be a complete metric space, let $S\subset X$, assume $V$ is a finite-dimensional Banach space, $L>0$ and $F\colon X \to V$ is an $L$-Lipschitz function. Let $\eta\in (0,1)$ and $\varepsilon>0$ and assume that $S$ is an $\tilde{A}(F,d)$ set. Finally, suppose $\mu$ is a Borel probability measure on $S$. Then there exists a function $G\colon X \to V$ and a compact set $E\subset S$ such that
    \begin{enumerate}
        \item\label{Enum:StV-main-1} the Lispchitz constant of $G$ satisfies $\tLip(G)\leq \tilde{K}(V,d) L$,
        \item\label{Enum:StV-main-2} for any $x\in X$,
        \begin{equation*}
            \lVert F(x) -G(x)\rVert_{V}\leq \varepsilon,
        \end{equation*}
        \item\label{Enum:StV-main-3} $\mu(E)\geq 1-\eta$ and $G(E)$ is contained in a finite union of $d$-dimensional affine subspaces of $V$.
    \end{enumerate}
\end{theorem}

\begin{lemma}\label{L:semicont}
    Suppose $S$ is a compact metric space, $Y$ is a metric space and $s\in (0,\infty)$. The functional
    \begin{equation*}
        f\mapsto \H^s_\infty(f(X))
    \end{equation*}
    is upper semi-continuous on the space $\tLip_L(X,Y)$ for any $L\in [0,\infty)$.
\end{lemma}

\begin{theorem}
    Suppose $X$ is a complete metric space and $S\subset X$.
    Suppose $d\in \nat$ and let $V$ be a finite-dimensional Banach space for which $\tilde{K}(V,d)=1$.
    Let $\mu$ be a $\sigma$-finite Borel measure on $S$ for which there exists a set $N\subset S$ with $\mu(N)=0$ such that $S\setminus N$ is $\sigma$-compact and an $\tilde{A}(F,d)$ set for every Lipschitz $F\colon X\to V$.
    Then, for any $L\in [0,\infty)$, the set
    \begin{equation*}
        \{f\in \tLip_L(X,V): \dim_H(f_\# \mu)\leq d\}
    \end{equation*}
    is residual and $G_\delta$ in $\tLip_L(X,V)$.
\end{theorem}

\begin{proof}
    We begin by proving the assertion assuming $S$ is a 
    compact $\tilde{A}(F,d)$ set and $\mu$ is a Borel probability measure on $S$, in which case we may simply take $N=\emptyset$. Let us write
    \begin{equation*}
        \mathcal{G}=\{f\in \tLip_L(X,V): \dim_H(f_\# \mu)\leq d\}
    \end{equation*}
    and
    \begin{equation*}
        \mathcal{G}_{\delta, \eta, \varepsilon}=\{f\in \tLip_L(X,V): \exists E \subset S\;\text{compact and such that,}\; \mu(E)\geq 1-\eta,\; \H^{d+\delta}_\infty(f(E)) < \varepsilon\},
    \end{equation*}
    for $\delta, \varepsilon>0$ and $\eta \in (0,\infty)$.
    Firstly, we shall assert that
    \begin{equation*}
        \mathcal{G}=\bigcap_{\varepsilon, \delta>0, \eta\in (0,1)} \mathcal{G}_{\delta, \eta, \varepsilon}.
    \end{equation*}
    The inclusion from left to right is immediate, we prove the converse inclusion. To that end, let $\delta>0$ and assume $f\in \mathcal{G}_{\delta, \eta, \varepsilon}$ for every $\eta\in (0,1)$ and $\varepsilon>0$. Then, by definition, for every $n,i,k\in \nat$, there exists a compact set $E(n,i,k)\subset S$ such that
    \begin{enumerate}
        \item\label{Enum:global-main-i} $\mu(S\setminus E(n,i,k))< \frac{1}{i}2^{-n}$,
        \item\label{Enum:global-main-ii} $\H^{d+\delta}_\infty(f(E(n,i,k)))< 2^{-i} \frac{1}{k}$.
    \end{enumerate}
    We let
    \begin{equation*}
        E_\delta=\bigcap_{k\in\nat} \bigcup_{i\in\nat} \bigcap_{n\in \nat} E(n,i,k).
    \end{equation*}
    Then \ref{Enum:global-main-i} and \ref{Enum:global-main-ii} imply 
    \begin{enumerate}[label={\rm(\alph*)}]
        \item $\mu(S\setminus E_\delta)=0$ and
        \item $\H^{d+\delta}_\infty(f(E_\delta))=0$.
    \end{enumerate}

    It follows that $\H^{d+\delta}(f(E_\delta))=0$, whence also $\dim_H(f(E_\delta))\leq d+\delta$. Taking
    \begin{equation*}
        E=\bigcap_{n\in \nat} E_{\frac{1}{n}},
    \end{equation*}
    yields $\mu(S\setminus E)=0$ and $\dim_H(f(E))\leq d$. In other words, $\dim_H(f_\# \mu)\leq d$ and hence $f\in \mathcal{G}$.

    Next, we show that for every $\delta, \varepsilon>0$, $\eta \in (0,1)$, the set $\mathcal{G}_{\delta, \eta, \varepsilon}$ is open and dense in $\tLip_L(X,V)$.
    Firstly, density. Fix $\delta>0$ and suppose $f\in \tLip_L(X,V)$. Since $S$ is an $\tilde{A}(f,d)$ set and $\tilde{K}(V,d)=1$, for any $\Delta>0$, we have from Theorem \ref{T:StV-main} (and Observation \ref{O:finiteness-subspaces}) a function $g\in \tLip_L(X,V)$ and a compact set $E\subset S$ such that $\mu(E)\geq 1-\eta$,
    \begin{equation*}
        \H^d(g(E))< \infty \quad \text{and}\quad \lVert f-g \rVert_{\infty}< \Delta.
    \end{equation*}
    This implies that $g\in \mathcal{G}_{\delta,\eta,\varepsilon}$ for every $\varepsilon>0$ and $\delta>0$. As $\Delta$ is arbitrary, this shows density of $\mathcal{G}_{\delta,\eta,\varepsilon}$ in $\tLip_L(X,V)$.

    To show openness, we write $\mathcal{K}(S)$ for the family of compact subsets of $S$ and
    \begin{equation*}
        \mathcal{G}_{\delta, \eta, \varepsilon}=\bigcup_{E\in \mathcal{K}(X); \mu(E)\geq 1-\eta}\{f\in \tLip_L(X,V): \H^{d+\delta}_\infty(f(E)) < \varepsilon\}.
    \end{equation*}
    Each of the sets inside the union on the right hand side is open due to Lemma \ref{L:semicont}. Therefore, the set $\mathcal{G}_{\delta, \eta, \varepsilon}$ is open.
    Since we may write
    \begin{equation*}
        \mathcal{G}=\bigcap_{\delta, \varepsilon>0, \eta\in(0,1)} \mathcal{G}_{\delta, \eta, \varepsilon}=\bigcap_{n,k,i\in \nat} \mathcal{G}_{\frac{1}{n}, \frac{1}{k}, \frac{1}{i}},
    \end{equation*}
    the set $\mathcal{G}$ is a countable intersection of open dense sets.

    In the general case, unless $\mu$ is trivial, in which case there is nothing to prove, it follows from  our assumptions, that there exists a set $ N\subset \tilde N \subset S$ and \emph{compact} $\tilde{A}(F,d)$ sets $S_1, S_2,\dots\subset S$ such that 
    \begin{enumerate}[label={\rm(\arabic*)}]
        \item $\mu(\tilde N)=0$,
        \item $0<\mu(S_i)<\infty$ and
        \item $S\setminus \tilde{N}= \bigcup_i S_i$.
    \end{enumerate}
    Denote $\mu_i= \frac{1}{\mu(S_i)}\mu_{|S_i}$. As each $S_i$ is compact and each $\mu_i$ is a probability measure on $S_i$, we have already shown that the sets
    \begin{equation*}
        \mathcal{F}_i=\{f\in \tLip_L(X,V): \dim_H(f_\# \mu_i)\leq d\}
    \end{equation*}
    are residual $G_\delta$ and thus, also the set
    \begin{equation*}
        \mathcal{F}=\bigcap_i \mathcal{F}_i
    \end{equation*}
    is residual $G_\delta$. It remains to show
    \begin{equation*}
        \{f\in \tLip_L(X,V): \dim_H(f_\# \mu)\leq d\}=\mathcal{F}.
    \end{equation*}
    The inclusion from left to right is immediate. Suppose $f\in \mathcal{F}$. Then, for each $i\in \nat$, there is a $E_i\subset S_i$ with $\mu(S_i\setminus E_i)=0$ and such that
    \begin{equation}\label{E:dim-leq-d}
        \dim_H(f(E_i))\leq d.
    \end{equation}
    Let $E=\bigcup_i E_i$. Then we have $f(E)=\bigcup_i f(E_i)$ and so by \eqref{E:dim-leq-d} we have
    \begin{equation*}
        \dim_H(f(E))= \dim_H(\bigcup f(E_i))\leq d,
    \end{equation*}
    showing that $\dim_H(f_\# \mu)\leq d$.
\end{proof}

\end{document}
